\documentclass[11pt]{book}

\usepackage[utf8]{inputenc}
\usepackage[T1]{fontenc}
\usepackage[most]{tcolorbox}
\usepackage{listings}
\usepackage{booktabs}
\usepackage{natbib}

\lstset{
  basicstyle=\ttfamily\small,
  columns=fullflexible,
  frame=single,
  breaklines=true,
}

\newtcolorbox{note}[1][]{colback=black!5,colframe=black!40,title={Note},#1}

\title{An Open-Source eBook}
\author{Noor Fatima}
\date{\today}

\begin{document}

\maketitle

\tableofcontents

\chapter{Getting Started}

\section{Why This Book Mixes Prose and Code}

A technical eBook reads best when code listings sit directly beside the
prose that explains them \citep{knuth1984texbook}. This starter uses
\texttt{listings} for that, kept plain rather than routed through a
shell-escape-dependent highlighter.

\begin{lstlisting}[caption={A minimal example.}]
def greet(name):
    return f"Hello, {name}!"
\end{lstlisting}

\begin{note}
Callout boxes like this one, built with \texttt{tcolorbox}, are useful for
asides that should stand apart from the main narrative.
\end{note}

\section{A Summary Table}

Table~\ref{tab:chapters} lists what each chapter of this starter covers.

\begin{table}[h]
  \centering
  \begin{tabular}{ll}
    \toprule
    Chapter & Covers \\
    \midrule
    1 & Code listings and callout boxes \\
    2 & Citations and the bibliography \\
    \bottomrule
  \end{tabular}
  \caption{Chapter overview.}
  \label{tab:chapters}
\end{table}

\chapter{Closing the Book}

\section{Citing a Source}

The reference used throughout this starter is resolved below
\citep{knuth1984texbook}.

\section{Wrapping Up}

The two chapters above are enough to show the combination this template is
for: prose, syntax-highlighted-looking listings, callouts, and a standard
bibliography, all under the plain \texttt{book} class.

\bibliographystyle{plainnat}
\bibliography{references}

\end{document}
